
%%%%%%%%%%%%%%%%%%%%%%%%%%%%%%%%%%%%%%%%%%%%%%%%%%%%%%%%%%%%%
\subsection{问题一的模型的建立和求解}
\subsubsection{具体分析}

问题一要求给出交会定位区域直径的计算算法，并回答以该直径为直径的圆能否覆盖定位区域，属于计算几何类问题。输入为若干检测点的平面坐标及其测得的示向度，输出为定位区域的凸多边形顶点、直径数值以及覆盖性的判定结论。由附录 2 可知，测向机在检测点处测得的是最大场强方向的方位角，因当地电磁环境的影响与真实方位角之间存在误差，该误差在目标区域全局范围内不超过 $1^\circ$，且同一地点的误差在一段时间内固定。因此单次观测给出的不是一条确定的直线，而是以检测点为顶点、张角为 $2^\circ$ 的楔形，多个检测点的楔形求交便得到干扰源的最可能位置集合。

本问的难点有三处。第一是角度区间的数值表达，示向度取值范围为 $[0^\circ,360^\circ)$，当读数接近 $0^\circ$ 或 $360^\circ$ 时直接按角度区间求交会因回绕而出错。第二是定位区域的有界性，两条视线近乎平行时楔形交会沿视线方向无限延展，此时直径无定义，必须引入目标区域约束加以截断。第三是覆盖判定的非平凡性，由 Jung 定理任意直径 $D_P$ 的平面点集其最小包围圆半径最大可达 $D_P/\sqrt3$，而直径圆的半径仅为 $D_P/2$，二者存在约 $15\%$ 的间隙，因此"以直径为直径的圆总能盖住原图形"的直觉并不成立，必须给出精确判据。本文的技术路线是：楔形半平面化，逐次半平面裁剪求交，凸包与旋转卡壳求直径，最后由 Thales 逆定理判定覆盖性并以大规模随机构型做统计验证。

\subsubsection{模型建立}

\paragraph{楔形的半平面等价表示} 记 $u(a)=(\cos a,\sin a)$，则位于 $S_i$、朝向 $\theta_i$、张角 $2\delta$ 的楔形为 $W_i=\{S_i+\rho u(a):\rho\ge0,\lvert a-\theta_i\rvert\le\delta\}$。注意到边界射线方位角为 $a$ 时其左法向量为 $n(a)=(-\sin a,\cos a)$，且对任意位移 $d$ 成立二维叉积恒等式
\begin{equation}
\operatorname{cross}\bigl(u(a),d\bigr)=d\cdot n(a),
\end{equation}
故 $d\cdot n(a)>0$ 当且仅当 $d$ 位于射线 $u(a)$ 的逆时针一侧。于是楔形可等价写为两个闭半平面的交
\begin{equation}
W_i=\bigl\{\,X:\ (X-S_i)\cdot n(\theta_i-\delta)\ge0,\ \ (X-S_i)\cdot n(\theta_i+\delta)\le0\,\bigr\}.
\label{eq:w-hp}
\end{equation}
式 \eqref{eq:w-hp} 不含任何角度区间的判断，从表达层面彻底消除了 $0^\circ$ 与 $360^\circ$ 处的回绕风险。

\paragraph{定位区域与截断} 干扰源须同时满足全部检测点的约束，故定位区域为各楔形与目标圆域的交
\begin{equation}
P=\bigcap_{i=1}^{m}W_i\ \cap\ D,\qquad D=\{X:\lVert X\rVert\le R\},\quad R=1800\ \text{m}.
\label{eq:region}
\end{equation}
$P$ 是有限个闭半平面的交，因而为闭凸集。算法上把式 \eqref{eq:region} 统一处理为：初始多边形取半宽 $M=20000$ 米的外接方框（远大于 $R$，用于区分"截断后的有界区域"与"未加截断时的无界退化构型"），再用 $2m$ 条楔形边界与目标圆域的内接 720 边形逐次裁剪。圆域用内接多边形近似是保守的，其相对收缩量为 $1-\cos(\pi/720)\approx9.5\times10^{-6}$，相对于千米量级的直径完全可忽略。

\paragraph{Sutherland--Hodgman 裁剪} 设当前凸多边形顶点序列为 $\{V_1,\dots,V_h\}$，用半平面 $\{X:(X-S)\cdot n\ge0\}$ 裁剪时依次考察每条有向边 $V_kV_{k+1}$：两端点均在可行侧则保留 $V_{k+1}$；均在外侧则全部舍弃；一内一外时保留交点
\begin{equation}
V_k+\lambda(V_{k+1}-V_k),\qquad
\lambda=\frac{(S-V_k)\cdot n}{(V_{k+1}-V_k)\cdot n}\in[0,1].
\label{eq:clip}
\end{equation}
该过程对每个半平面耗时 $O(h)$，总复杂度为 $O(mh)$。

\paragraph{直径算法} 凸多边形 $P$ 的直径定义为 $D_P=\max_{a,b\in P}\lVert a-b\rVert$。由于欧氏范数为凸函数而 $P$ 为凸紧集，最大值必在顶点上取到，故 $D_P=\max_{j<k}\lVert V_j-V_k\rVert$。本文实现两种算法互为校验：其一是顶点对暴力枚举，复杂度 $O(h^2)$，直接且无附加假设；其二是旋转卡壳法，先以 Andrew 单调链求凸包以去除裁剪产生的共线冗余点，再沿边界同步推进两条对踵切线，复杂度 $O(h)$。

\paragraph{直径圆覆盖的充要判据} 设定位区域的直径端点为 $A,C$，以 $AC$ 为直径的圆 $\Gamma$ 的圆心为 $(A+C)/2$、半径为 $D_P/2$。由 Thales 定理的逆定理，点 $X$ 落在 $\Gamma$ 内当且仅当 $(X-A)\cdot(X-C)\le0$，即 $\angle AXC\ge90^\circ$。因此对凸四边形 $ABCD$，若 $\lVert AC\rVert=D_P$，则被 $\Gamma$ 覆盖当且仅当
\begin{equation}
\boxed{\ \angle ABC\ge90^\circ\ \ \text{且}\ \ \angle ADC\ge90^\circ\ }
\label{eq:thales}
\end{equation}
等价地可用符号量 $\sigma_B=(B-A)\cdot(B-C)$ 与 $\sigma_D=(D-A)\cdot(D-C)$ 表述。由四边形内角和为 $360^\circ$ 可知，判据 \eqref{eq:thales} 成立当且仅当 $\angle A+\angle C\le180^\circ$，这正是"瘦长四边形易被覆盖"这一直觉的精确表述。作为对照，Jung 定理给出最小包围圆半径的界 $\rho_{\min}\le D_P/\sqrt3$，与 $D_P/2$ 之比为 $\sqrt3/2\approx0.866$，说明覆盖失败在一般构型下确实可能发生。

\subsubsection{模型求解}

取示范构型：检测点 $S_1=(0,0)$、$S_2=(1200,0)$，真实干扰源位于 $G=(900,700)$，真实方位角为 $37.875^\circ$ 与 $113.199^\circ$，叠加环境偏差后示向度读数为 $38.475^\circ$ 与 $112.399^\circ$。求解得四顶点凸四边形，直径 $D_P=56.7078$ 米，旋转卡壳与暴力枚举的偏差为零，最小包围圆半径为 $28.3539$ 米恰等于 $D_P/2$，即直径圆恰好外接于该区域；两个 Thales 符号量分别为 $-332.84$ 与 $-323.15$，均为负，覆盖成立。如图 \ref{fig:q1-region} 所示，两条示向度射线束在目标附近相夹形成狭长的凸四边形，其长对角线的中点恰为直径圆圆心。

\begin{figure}[htbp]
	\centering
	\includegraphics[width=0.8\textwidth]{../求解/问题一/figures/交会定位区域示意图.png}
	\caption{交会定位区域、直径与直径圆（示范构型）}
	\label{fig:q1-region}
\end{figure}

为回答覆盖性的普遍问题，本文在目标区域内随机生成检测点与干扰源，以真实方位角叠加 $U[-1^\circ,1^\circ]$ 的固定偏差构造示向度，得到 $5000$ 组构型。由于示向度由真实位置反推，真值必然落在定位区域内，构型始终具有物理意义。统计采用两种口径：纯交会四边形（题目图 2 所示的红色四边形原型）与叠加目标圆域截断后的定位区域。

\begin{table}[H]
	\centering
	\caption{两种口径下的蒙特卡洛统计（各 5000 组构型）}
	\label{tab:q1-mc}
	\begin{tabularx}{\textwidth}{LCC}
		\toprule
		统计量 & 纯交会四边形 & 加目标圆域截断 \\
		\midrule
		非空构型占比 & $100.00\%$ & $100.00\%$ \\
		无界退化构型占比 & $2.80\%$ & $2.80\%$ \\
		有效构型数 & $4860$ & $4860$ \\
		直径圆覆盖率 & $\mathbf{98.99\%}$ & $96.28\%$ \\
		直径均值 (m) & $402.53$ & $271.04$ \\
		直径中位数 (m) & $157.96$ & $151.88$ \\
		最小包围圆半径与直径之比（均值） & $0.500001$ & $0.5002$ \\
		最小包围圆半径与直径之比（最大） & $0.500280$ & $0.5611$ \\
		\bottomrule
	\end{tabularx}
\end{table}

由表 \ref{tab:q1-mc} 可见，纯交会口径下直径圆覆盖率为 $98.99\%$；未加截断时约有 $2.8\%$ 的构型因两条视线近乎平行而退化无界，这部分构型在计入目标区域约束后自动变为有界，因此截断是必需的。截断口径下覆盖率降至 $96.28\%$，原因是定位区域与圆域边界相交后成为含圆弧边的多边形，圆弧段必然外凸于任何以其弦为直径的圆。图 \ref{fig:q1-cover} 对比了覆盖成立与覆盖失败的两个真实算例：左图中定位区域近似为平行四边形，四个顶点全部落在直径圆内；右图取自蒙特卡洛中超出量最大的构型，其右侧顶点明显突出于直径圆之外，Thales 符号量一正一负。

\begin{figure}[htbp]
	\centering
	\includegraphics[width=0.8\textwidth]{../求解/问题一/figures/直径圆覆盖与不覆盖对比图.png}
	\caption{直径圆覆盖交会定位区域的两种情形}
	\label{fig:q1-cover}
\end{figure}

按两视线在干扰源处的交会角 $\alpha$ 分组统计，结果如表 \ref{tab:q1-group} 所示。

\begin{table}[H]
	\centering
	\caption{直径圆覆盖率按交会角分组统计（纯交会口径）}
	\label{tab:q1-group}
	\begin{tabular}{cccccc}
		\toprule
		交会角区间 & 构型数 & 直径均值 (m) & 直径中位数 (m) & 覆盖率 & 包围圆半径/直径均值 \\
		\midrule
		$(0^\circ,30^\circ]$    & $1486$ & $950.26$ & $499.79$ & $100.00\%$ & $0.5000$ \\
		$(30^\circ,45^\circ]$   & $688$  & $201.55$ & $200.85$ & $100.00\%$ & $0.5000$ \\
		$(45^\circ,60^\circ]$   & $568$  & $136.17$ & $134.21$ & $100.00\%$ & $0.5000$ \\
		$(60^\circ,75^\circ]$   & $482$  & $106.85$ & $107.16$ & $100.00\%$ & $0.5000$ \\
		$(75^\circ,90^\circ]$   & $404$  & $83.50$  & $85.25$  & $98.02\%$  & $0.5000$ \\
		$(90^\circ,105^\circ]$  & $342$  & $76.21$  & $77.74$  & $88.01\%$  & $0.5000$ \\
		$(105^\circ,120^\circ]$ & $231$  & $84.82$  & $87.31$  & $100.00\%$ & $0.5000$ \\
		$(120^\circ,150^\circ]$ & $363$  & $111.35$ & $112.21$ & $100.00\%$ & $0.5000$ \\
		$(150^\circ,180^\circ]$ & $296$  & $530.07$ & $299.21$ & $100.00\%$ & $0.5000$ \\
		\bottomrule
	\end{tabular}
\end{table}

表 \ref{tab:q1-group} 揭示了一个清晰的结构：覆盖率仅在交会角接近 $90^\circ$ 的区间出现明显下降，其余区间均为 $100\%$。其几何原因在于，当两视线正交时两个张角 $2^\circ$ 的锥相交所得的四边形趋近于矩形，而矩形恰是 Thales 判据的临界形状——四个内角恰为 $90^\circ$、四个顶点恰好落在直径圆上；真实锥形带来的微小扰动使其中某个顶点稍微外移或内移，于是覆盖与否各占一定比例。当交会角显著偏离 $90^\circ$ 时，定位区域变为长菱形，由等价条件 $\angle A+\angle C\le180^\circ$ 可知其必然满足覆盖判据。此外交会角越小则定位区域直径越大，$(0^\circ,30^\circ]$ 区间的直径均值接近 $1$ 千米，这正是交会几何变差的直接体现。

\begin{figure}[htbp]
	\centering
	\includegraphics[width=0.8\textwidth]{../求解/问题一/figures/直径圆覆盖率与直径分布图.png}
	\caption{定位区域直径分布、覆盖率随交会角的变化与测向精度灵敏度}
	\label{fig:q1-stat}
\end{figure}

如图 \ref{fig:q1-stat}(a) 所示，定位区域直径呈显著右偏分布，中位数约 $158$ 米而均值超过 $400$ 米，说明大部分构型下定位精度尚可，但存在少量直径达千米量级的劣质构型。图 \ref{fig:q1-stat}(b) 直观呈现了覆盖率在 $90^\circ$ 附近的凹陷。图 \ref{fig:q1-stat}(c) 给出覆盖率对测向精度上界 $\delta$ 的灵敏度：当 $\delta$ 由 $1^\circ$ 收紧至 $0.25^\circ$ 时覆盖率升至 $99.93\%$，放宽至 $4^\circ$ 时降至 $96.39\%$，同时失败构型的最大超出量由 $1.17$ 米增至 $19.45$ 米。需要强调的是，全部失败构型中顶点超出直径圆的最大距离仅为 $1.1697$ 米，平均仅 $0.4928$ 米，占定位区域直径的比例不超过 $1.37\%$。

综上，以定位区域直径为直径的圆不能保证覆盖定位区域，其成立的充要条件是式 \eqref{eq:thales}；在满足有界性的随机交会构型中覆盖成立的频率约为 $98.99\%$，失效仅发生在交会角接近 $90^\circ$ 的临界构型上，且失效的绝对幅度不超过米级，对后续以 $20$ 米清除半径为精度要求的任务而言影响可以忽略。
